\newcommand{\titulus}{\nomenFesti{Dominica XVIII per Annum.}}
\newcommand{\tedeumsimplex}{Simplex}
\newcommand{\oratio}{\pars{Oratio.}

\noindent Adésto, Dómine, fámulis tuis et perpétuam benignitátem largíre poscéntibus, ut his, qui te auctórem et gubernatórem gloriántur habére, et grata restáures et restauráta consérves.

\pars{Pro pace in Ucraina.} \scriptura{Sir. 50, 25; 2 Esdr. 4, 20; \textbf{H416}}

\vspace{-4mm}

\antiphona{II D}{temporalia/ant-dapacemdomine.gtex}

\vfill

\noindent Deus, a quo sancta desidéria, recta consília et iusta sunt ópera: da servis tuis illam, quam mundus dare non potest, pacem; ut et corda nostra mandátis tuis dédita, et hóstium subláta formídine, témpora sint tua protectióne tranquílla.

\noindent Per Dóminum nostrum Iesum Christum, Fílium tuum, qui tecum vivit et regnat in unitáte Spíritus Sancti, Deus, per ómnia sǽcula sæculórum.

\noindent \Rbardot{} Amen.}
\newcommand{\hymnusmatutinum}{\pars{Hymnus.}

\vspace{-5mm}

\antiphona{II}{temporalia/hym-MediaeNoctis.gtex}}
\newcommand{\nocturnoii}{\vspace{-4mm}

\pars{Psalmus 4.}

\vspace{-4mm}

\antiphona{III a}{temporalia/ant-confessionemetdecorem.gtex}

%\vspace{-2mm}

\scriptura{Ps. 103, 1-12}

%\vspace{-2mm}

\initiumpsalmi{temporalia/ps103i-initium-iii-a-auto.gtex}

\input{temporalia/ps103i-iii-a.tex} \Abardot{}

\vfill
\pagebreak

\pars{Psalmus 5.}

\vspace{-4mm}

\antiphona{II D}{temporalia/ant-dominumdeumadoremus.gtex}

%\vspace{-2mm}

\scriptura{Ps. 103, 13-23}

\initiumpsalmi{temporalia/ps103ii-initium-ii-D-auto.gtex}

\input{temporalia/ps103ii-ii-D.tex} \Abardot{}

\vfill
\pagebreak

\pars{Psalmus 6.} \scriptura{Ps. 103, 24}

\vspace{-4mm}

\antiphona{E}{temporalia/ant-quammagnificatasunt.gtex}

\vspace{-4mm}

\scriptura{Ps. 103, 24-35}

%\vspace{-2mm}

\initiumpsalmi{temporalia/ps103iii-initium-e.gtex}

\vspace{-1.5mm}

\input{temporalia/ps103iii-e.tex} \Abardot{}

\vfill
\pagebreak}
\newcommand{\nocturnoiii}{\pars{Cantica.}

\vspace{-4mm}

\antiphona{D}{temporalia/ant-eccedeusnoster.gtex}

%\vspace{-2mm}

\scriptura{Canticum Isaiæ, Is. 33, 2-10}

%\vspace{-2mm}

\initiumpsalmi{temporalia/isaiae7-initium-d-g-auto.gtex}

\input{temporalia/isaiae7-d-g.tex} \hfill \rubrica{Hic non dicitur antiphona.}

\vfill
\pagebreak

\scriptura{Canticum Isaiæ, Is. 33, 13-17}

%\vspace{-2mm}

\initiumpsalmi{temporalia/isaiae8-initium-d-g-auto.gtex}

\input{temporalia/isaiae8-d-g.tex}

\vfill
\pagebreak

\scriptura{Canticum Ecclesiastici, Sir. 36, 14-19}

%\vspace{-2mm}

\initiumpsalmi{temporalia/ecclesiasticus36-initium-d-g-auto.gtex}

\input{temporalia/ecclesiasticus36-d-g.tex}

\vfill

\antiphona{}{temporalia/ant-eccedeusnoster.gtex}

\vfill
\pagebreak}
\newcommand{\lectioi}{\pars{Lectio I.} \scriptura{Prov. 1, 1-7}

\noindent Incipit liber Proverbiórum.

\noindent Parábolæ Salomónis fílii David regis Israel

\noindent ad sciéndam sapiéntiam et disciplínam,

\noindent ad intellegénda verba prudéntiæ;

\noindent ad suscipiéndam eruditiónem doctrínæ, iustítiam et iudícium et æquitátem,

\noindent ut detur párvulis astútia, adulescénti sciéntia et recogitátio.

\noindent Audiat sápiens et addet doctrínam, et intéllegens dispositiónes possidébit:

\noindent animadvértet parábolam et allegoríam, verba sapiéntium et ænígmata eórum.

\noindent Timor Dómini princípium sciéntiæ. Sapiéntiam atque doctrínam stulti despíciunt.}
\newcommand{\responsoriumi}{\pars{Responsorium 1.} \scriptura{\Rbardot{} Ps. 110, 10 \Vbardot{} Sap. 6, 19 \& Cantor; \textbf{H399}}

\vspace{-5mm}

\responsorium{II}{temporalia/resp-initiumsapientiae-CROCHU.gtex}{}}
\newcommand{\lectioii}{\pars{Lectio II.} \scriptura{Prov. 1, 8-19}

\noindent Audi, fili mi, disciplínam patris tui,

\noindent et ne dimíttas legem matris tuæ:

\noindent ut addátur grátia cápiti tuo, et torques collo tuo.

\noindent Fili mi, si te lactáverint peccatóres, ne acquiéscas eis.

\noindent Si díxerint: Veni nobíscum, insidiémur sánguini;

\noindent abscondámus tendículas contra insóntem frustra;

\noindent deglutiámus eum sicut inférnus vivéntem, et íntegrum quasi descendéntem in lacum;

\noindent omnem pretiósam substántiam reperiémus;

\noindent implébimus domos nostras spóliis:

\noindent sortem mitte nobíscum, marsúpium unum sit ómnium nostrum:

\noindent fili mi, ne ámbules cum eis; próhibe pedem tuum a sémitis eórum:

\noindent pedes enim illórum ad malum currunt, et festínant ut effúndant sánguinem.

\noindent Frustra autem iácitur rete ante óculos pennatórum.

\noindent Ipsi quoque contra sánguinem suum insidiántur, et moliúntur fraudes contra ánimas suas.

\noindent Sic sémitæ omnis avári: ánimas possidéntium rápiunt.}
\newcommand{\responsoriumii}{\pars{Responsorium 2.} \scriptura{\Rbardot{} Prov. 1, 8 \Vbardot{} ibid. 3, 9; \textbf{H401}}

\vspace{-5mm}

\responsorium{VII}{temporalia/resp-audifilimihidisciplinam-CROCHU.gtex}{}}
\newcommand{\lectioiii}{\pars{Lectio III.} \scriptura{Prov. 1, 20-33}

\noindent Sapiéntia foris prǽdicat, in platéis dat vocem suam,

\noindent in cápite viárum frequéntium clámitat,

\noindent in fóribus portárum urbis profert verba sua:

\noindent «Usquequo, párvuli, dilígitis infántiam,

\noindent et derisóres sibi derisiónem cúpient,

\noindent et imprudéntes odíbunt sciéntiam?

\noindent Convertímini ad correptiónem meam;

\noindent en próferam vobis spíritum meum

\noindent et osténdam vobis verba mea.

\noindent Quia vocávi, et renuístis,

\noindent exténdi manum meam, et non fuit qui aspíceret;

\noindent despexístis omne consílium meum

\noindent et increpatiónes meas neglexístis.

\noindent Ego quoque in intéritu vestro ridébo

\noindent et subsannábo, cum terror vobis advénerit,

\noindent cum irrúerit ut procélla terror,

\noindent et intéritus quasi tempéstas ingrúerit,

\noindent quando vénerit super vos tribulátio et angústia».

\noindent Tunc invocábunt me, et non exáudiam,

\noindent instánter quærent me, et non invénient me,

\noindent eo quod exósam habúerint disciplínam

\noindent et timórem Dómini non elégerint,

\noindent nec acquiéverint consílio meo

\noindent et despéxerint univérsam correptiónem meam.

\noindent Cómedent ígitur fructus viæ suæ

\noindent suísque consíliis saturabúntur.

\noindent Avérsio parvulórum interfíciet eos,

\noindent et secúritas stultórum perdet illos.

\noindent Qui autem me audíerit, absque terróre requiéscet,

\noindent et tranquíllus erit timóre malórum subláto.}
\newcommand{\responsoriumiii}{\pars{Responsorium 3.} \scriptura{\Rbardot{} Prov. 1, 32-33 \Vbardot{} ibid. 8, 4; \textbf{H401}}

\vspace{-5mm}

\responsorium{VI}{temporalia/resp-aversioparvulorum-CROCHU-cumdox.gtex}{}}
\newcommand{\lectioiv}{\pars{Lectio IV.} \scriptura{1, 18-19: SCh 121, 52-53}

\noindent Ex Commentário sancti Ephræm diáconi in Diatéssaron.

\noindent Quis unum ex effátis tuis, Dómine, mente penetráre póterit?

\noindent Plus est quod relínquimus quam quod cápimus, sicut sitiéntes qui bibunt ex fonte. Nam verbum Dómini iuxta multas discéntium perceptiónes multos præbet aspéctus. Dóminus verbum suum multis colóribus depínxit, ut, quisquis dísceret, in id inspíceret quod ei placéret. Vários thesáuros in verbo suo cóndidit, ut quisque nostrum, ubi se exercéret, inde ditésceret.

\noindent Verbum Dei arbor vitæ est, quæ ex ómnibus suis pártibus fructum benedíctum tibi offert, sicut illa rupes quæ in desérto aperiebátur, ut ómnibus suis pártibus potum præbéret spiritálem. \emph{Edébant}, ait Apóstolus, \emph{spiritálem cibum, et bibébant potum spiritálem.}}
\newcommand{\responsoriumiv}{\pars{Responsorium 4.} \scriptura{\Rbardot{} Prv. 8, 23-25 \Vbardot{} ibid. 8, 26.30; \textbf{H398}}

\vspace{-5mm}

\responsorium{I}{temporalia/resp-inprincipiodeusantequam-CROCHU.gtex}{}}
\newcommand{\lectiov}{\pars{Lectio V.}

\noindent Cui ergo áliqua pars ex thesáuro eius contíngit, ne credat id solum, quod ipse invénit, in hoc verbo inésse, sed exístimet se id solum ex multis quæ in eo sunt, inveníre potuísse.

\noindent Nec proptérea quod hæc pars sola ad eum pérvenit eíque contíngit, ipsum verbum exíle et stérile dicat atque despíciat, sed, quia id cápere non potest, propter divítias eius grátias agat.

\noindent Gaude quod victus es, neque contristéris quod te superávit.

\noindent Sítiens gaudet cum bibit, nec contristátur quod fontem exhauríre non potest.

\noindent Vincat fons sitim tuam, non autem sitis fontem vincat, quia, si sitis tua explétur quin fons exhauriátur, dénuo sítiens íterum ex eo bíbere póteris; si vero, siti tua expléta, fons quoque siccarétur, victória tua in malum tuum verterétur.}
\newcommand{\responsoriumv}{\pars{Responsorium 5.} \scriptura{\Rbardot{} Sap. 9, 10 \Vbardot{} ibid. 9, 4; \textbf{H398}}

\vspace{-5mm}

\responsorium{I}{temporalia/resp-emittedominesapientiam-CROCHU-sinedox.gtex}{}}
\newcommand{\lectiovi}{\pars{Lectio VI.}

\noindent Grátias age pro eo quod accepísti, et propter id quod remánsit et abundávit noli contristári.

\noindent Quod accepísti et ad quod pervenísti, pars tua est, et id quod remánsit, tua est heréditas.

\noindent Quod propter infirmitátem tuam una hora accípere non potes, áliis horis, si perseveráveris, accípere póteris.

\noindent Nec mente malígna conéris uno haustu súmere, quod uno haustu sumi nequit, neque ex ignávia desístas ab eo quod paulátim súmere possis.}
\newcommand{\responsoriumvi}{\pars{Responsorium 6.} \scriptura{\Rbardot{} Sap. 9, 4.5 \Vbardot{} ibid. 9, 10; \textbf{H399}}

\vspace{-5mm}

\responsorium{I}{temporalia/resp-damihidominesedium-CROCHU-cumdox.gtex}{}}
\newcommand{\evangelium}{
\pars{Versus.} \scriptura{Ps. 118, 148}

% Versus. %%%
\sineinitiali{temporalia/versus-praevenerunt.gtex}

\vspace{5mm}

\sineinitiali{temporalia/oratiodominica-mat.gtex}

\vspace{5mm}

\pars{Absolutio.}

\cuminitiali{}{temporalia/absolutio-avinculis.gtex}

\vfill
\pagebreak

\cuminitiali{}{temporalia/benedictio-solemn-evangelica.gtex}

\vspace{7mm}

\pars{Evangelium} \scriptura{Mt. 14, 13-21}

\noindent Léctio sancti Evangélii secúndum Matthǽum.

\noindent In illo témpore: cum audísset Iesus [mortem Ioánnis Baptístæ], secéssit inde in navícula in locum desértum seórsum; et cum audíssent, turbæ secútæ sunt eum pedéstres de civitátibus. Et éxiens vidit turbam multam et misértus est eórum et curávit lánguidos eórum.

\noindent Véspere autem facto, accessérunt ad eum discípuli dicéntes: «Desértus est locus, et hora iam prætériit; dimítte turbas, ut eúntes in castélla emant sibi escas».

\noindent Iesus autem dixit eis: «Non habent necésse ire; date illis vos manducáre».

\noindent Illi autem dicunt ei: «Non habémus hic nisi quinque panes et duos pisces».

\noindent Qui ait: «Afférte illos mihi huc».

\noindent Et cum iussísset turbas discúmbere supra fenum, accéptis quinque pánibus et duóbus píscibus, aspíciens in cælum benedíxit et fregit et dedit discípulis panes, discípuli autem turbis.

\noindent Et manducavérunt omnes et saturáti sunt; et tulérunt relíquias fragmentórum duódecim cóphinos plenos. Manducántium autem fuit númerus fere quinque mília virórum, excéptis muliéribus et párvulis.

\scriptura{Lib. 2,3: SC 93,252-254}

\noindent Ex Tractátu Balduíni Cantuariénsis epíscopi De sacraménto altáris.

\noindent \emph{Qui edunt me adhuc esúrient, et qui bibunt me adhuc sítient.}

\noindent Christus vero, qui est Dei sapiéntia, in præsénti non manducátur in satietáte desidérii, sed in desidério satietátis: et quanto plus eius suávitas gustátur, tanto magis desidérium irritátur. Proptérea qui edunt adhuc esúrient, donec véniat satíetas. Cum autem in bonis replétum fúerit desidérium, tunc neque esúrient neque sítient ámplius.

\noindent Potest étiam secúndum futúrum statum intéllegi quod dictum est: \emph{Qui edunt me adhuc esúrient}. Est enim in illa ætérna satietáte quasi quædam esúries, non egestátis sed felicitátis, ubi semper manducáre áppetunt, qui manducáre numquam volunt, et in satietáte numquam fastídiunt. Est enim satíetas sine fastídio et desidérium sine suspírio.

\noindent {\color{gray} Christus enim in decóre suo semper admirábilis, semper est et desiderábilis, \emph{in quem desíderant ángeli prospícere}.

\noindent Ergo et cum habétur desiderátur, et cum tenétur quǽritur, sicut scriptum est: \emph{Quǽrite fáciem eius semper}. Semper enim quǽritur, qui ad hoc dilígitur ut semper habeátur. Itaque et qui invéniunt adhuc quærunt, et qui edunt adhuc esúriunt, et qui bibunt adhuc sítiunt.

\noindent Sed hæc inquisítio omnem sollicitúdinem tollit; sed hæc esúries omnem esúriem pellit; sed hæc sitis omnem sitim exstínguit. Non est hæc egestátis sed consummátæ felicitátis.

\noindent De illa enim quæ egestátis est dícitur: \emph{Qui venit ad me non esúriet, et qui credit in me non sítiet umquam.}

\noindent De ea vero quæ felicitátis est dícitur: \emph{Qui edunt me adhuc esúrient, et qui bibunt me adhuc sítient.}

\noindent Christus in se credéntibus cibus est et potus, panis et vinum: panis cum confírmat et sólidat, potus vel vinum est cum lætíficat.

\noindent Quod si Christus secúndum præséntem iustórum fortitúdinem et lætítiam panis est et potus, quanto magis in futúro hoc erit secúndum quod iustis tunc erit?}

\vfill
\pagebreak

\pars{Responsorium 7.} \scriptura{\Rbardot{} Io. 6, 48 \Vbardot{} ibid. 6, 51.52}

\vspace{-5mm}

\responsorium{VII}{temporalia/resp-egosumpanisvitae-E611-cumdox.gtex}{}

\vfill
\pagebreak
}
\newcommand{\benedictus}{\pars{Canticum Zachariæ.} \scriptura{Cf. Mc. 8, 6; \textbf{H429}}

%\vspace{-4mm}

{
\grechangedim{interwordspacetext}{0.18 cm plus 0.15 cm minus 0.05 cm}{scalable}%
\antiphona{III a\textsuperscript{2}}{temporalia/ant-accipiensdominusquinque.gtex}
\grechangedim{interwordspacetext}{0.22 cm plus 0.15 cm minus 0.05 cm}{scalable}%
}

\vspace{-3mm}

\scriptura{Lc. 1, 68-79}

\vspace{-2mm}

\initiumpsalmi{temporalia/benedictus-initium-iiisoll-a4-auto.gtex}

\vspace{-1.5mm}

\input{temporalia/benedictus-iiisoll-a4.tex}

{
\grechangedim{interwordspacetext}{0.18 cm plus 0.15 cm minus 0.05 cm}{scalable}%
\antiphona{}{temporalia/ant-accipiensdominusquinque.gtex}
\grechangedim{interwordspacetext}{0.22 cm plus 0.15 cm minus 0.05 cm}{scalable}%
}
}
\include{hebdomadaxviii}
\include{dominica}
